\begin{table}[H]
    \centering
    \caption{唤醒源}
    \label{tab:lpm-wakeup-source}
    \begin{threeparttable}
    \begin{tabular}{| r | l | c | c | c |}
    \hline
    \rowcolor{lightgray}
    \textbf{WAKEUP\_ENA} & \textbf{唤醒源\tnote{1}} &  \textbf{Light-sleep} & \textbf{Deep-sleep} & \textbf{休眠} \\\hline
        0x1 & EXT0\tnote{2} & Y & Y & - \\\hline
        0x2 & EXT1\tnote{3} & Y & Y & Y \\\hline
        0x4 & GPIO\tnote{4} & Y & Y & - \\\hline
        0x8 & RTC 计时器 & Y & Y & Y \\\hline
        0x10 & SDIO\tnote{5} & Y & - & - \\\hline
        0x20 & Wi-Fi\tnote{6} & Y & - & - \\\hline
        0x40 & UART0 \tnote{7}& Y & - & - \\\hline
        0x80 & UART1 \tnote{7} & Y & - & - \\\hline
        0x100 & TOUCH & Y & Y & - \\\hline
        0x200 & ULP 协处理器 & Y & Y & - \\\hline
        0x400 & BT \tnote{6}& Y & - & -  \\\hline
    \end{tabular}
    \begin{tablenotes}
        \item[1] 表格中只有 GPIO 和 SDIO 可以配置为拒绝睡眠原因。
        \item[2] EXT0 仅可将芯片从 Light-sleep/Deep-sleep 模式中唤醒。当 \hyperref[regdesc:RTCCNTLEXTWAKEUPCONFREG]{RTC\_CNTL\_EXT\_WAKEUP0\_LV} 为 1 时将触发管脚高电平，否则将触发管脚低电平。用户可通过设置 \hyperref[regdesc:RTCIOEXTWAKEUP0REG]{RTCIO\_EXT\_WAKEUP0\_SEL[4:0]}，选择将作为唤醒源的 RTC 管脚。
        \item[3] EXT1 经过专门设计，可将芯片从任何睡眠模式中唤醒，而且还支持多个管脚的组合。 首先，应按照选定作为唤醒源的管脚位图，配置 \hyperref[regdesc:RTCCNTLEXTWAKEUP1REG]{RTC\_CNTL\_EXT\_WAKEUP1\_SEL[17:0]}。接着，当 \hyperref[regdesc:RTCCNTLEXTWAKEUPCONFREG]{RTC\_CNTL\_EXT\_WAKEUP1\_LV} 为 1 时，只要有任何选择的管脚为高电压，则将触发信号唤醒芯片；当 \hyperref[regdesc:RTCCNTLEXTWAKEUPCONFREG]{RTC\_CNTL\_EXT\_WAKEUP1\_LV} 为 0 时，则必须全部选择的管脚均为低电压才会触发信号。

        \textbf{注意：}用户应确保 EXT1 信号的 hold 时长至少三个RTC\_SLOW\_CLK 周期，否则信号状态无法被 \hyperref[regdesc:RTCCNTLEXTWAKEUP1STATUSREG]{RTC\_CNTL\_EXT\_WAKEUP1\_STATUS} 记录。
        \item[4] 在 Deep-sleep 模式下，仅有 RTC GPIO 可以作为唤醒源，而非数字 GPIO。
        \item[5] 任何接收到的 SDIO 命令均将触发唤醒动作。
        \item[6] 为了通过 Wi-Fi 或 BT 源唤醒芯片，芯片将在 Active、Modem-sleep 和 Light-sleep 之间进行切换，CPU、Wi-Fi、Bluetooth 和射频模块均将在预设间隔中唤醒，保证 Wi-Fi/蓝牙的正常连接。
        \item[7] 当输入 RxD 沿变化的次数大于等于 (UART\_ACTIVE\_THRESHOLD+2) 时，即触发唤醒。{\bfseries 注意：}RxD 不能通过 GPIO 交换矩阵输入，只能通过 IO\_MUX 输入。
    \end{tablenotes}
    \end{threeparttable}
\end{table}
